\documentclass[a4paper,10pt]{report}\input{../0.Préambule/Préambule_cours_prof}\begin{document}

\pagestyle{empty}
\newpage
\section*{Agrandissement et réduction}
\addcontentsline{toc}{section}{Agrandissement et réduction}


\begin{exer}
Le triangle $ABC$ est un agrandissement de $IJK$. \\ Calculer $AB$ et $AC$.
\end{exer}

\begin{exer}
Tracer un agrandissement du parallélogramme \\$KLMN$ tel que $KL = 2$cm ,$LM = 3$cm et $\widehat{KLM} = 40$\degre sachant \\que le côté correspondant à $KL$ mesure alors $4,6$cm.
\end{exer}

\begin{exer}
Un triangle $A'B'C'$ rectangle en $A'$ et d'aire $27$cm$^2$ \\est un agrandissement d'un triangle $ABC$, rectangle en $A$ tel que \\$AB=3$cm et $AC=2$cm. Calcule les longueurs $A'B'$ et $A'C'$.
\end{exer}

\vspace{-5.8cm}
\begin{flushright}
\includegraphics[scale=0.3]{../40.Image/15. Act - Agrandissement et réduction 1.png}
\end{flushright}

\vspace{-7mm}
\begin{exer}

\medskip
\begin{enumerate}
\item Un rectangle a pour dimensions $15$cm et $12$cm. Calculer son périmètre.
\item Quel coefficient de réduction choisir pour obtenir un rectangle de $36$cm de périmètre ?
Justifier
\end{enumerate}
\end{exer}

\begin{exer}
La pyramide du Louvre est un pyramide régulière à base carrée de coté $35$m  et de hauteur $22$m.
\begin{enumerate}
\item Fais un schéma
\item Calcule le volume $\mathcal{V}$ de cette pyramide. Donne la valeur exacte en m$^3$ puis la valeur arrondie à l'unité.
\item Sur une maquette, on construit une réduction de cette pyramide, le coté de la base carrée mesure $7$cm. Calcule le coefficient de réduction.
\item Déduis-en le volume $\mathcal{V'}$ de la pyramide sur la maquette. donne la valeur exacte en cm$^3$ puis la valeur arrondie à l'unité.
\end{enumerate}
\end{exer}

\vspace{4cm}
\setcounter{exercice}{0}
\section*{Agrandissement et réduction}
\addcontentsline{toc}{section}{Agrandissement et réduction}


\begin{exer}
Le triangle $ABC$ est un agrandissement de $IJK$. \\ Calculer $AB$ et $AC$.
\end{exer}

\begin{exer}
Tracer un agrandissement du parallélogramme \\$KLMN$ tel que $KL = 2$cm ,$LM = 3$cm et $\widehat{KLM} = 40$\degre sachant \\que le côté correspondant à $KL$ mesure alors $4,6$cm.
\end{exer}

\begin{exer}
Un triangle $A'B'C'$ rectangle en $A'$ et d'aire $27$cm$^2$ \\est un agrandissement d'un triangle $ABC$, rectangle en $A$ tel que \\$AB=3$cm et $AC=2$cm. Calcule les longueurs $A'B'$ et $A'C'$.
\end{exer}

\vspace{-5.8cm}
\begin{flushright}
\includegraphics[scale=0.3]{../40.Image/15. Act - Agrandissement et réduction 1.png}
\end{flushright}

\vspace{-7mm}
\begin{exer}

\medskip
\begin{enumerate}
\item Un rectangle a pour dimensions $15$cm et $12$cm. Calculer son périmètre.
\item Quel coefficient de réduction choisir pour obtenir un rectangle de $36$cm de périmètre ?
Justifier
\end{enumerate}
\end{exer}

\begin{exer}
La pyramide du Louvre est un pyramide régulière à base carrée de coté $35$m  et de hauteur $22$m.
\begin{enumerate}
\item Fais un schéma
\item Calcule le volume $\mathcal{V}$ de cette pyramide. Donne la valeur exacte en m$^3$ puis la valeur arrondie à l'unité.
\item Sur une maquette, on construit une réduction de cette pyramide, le coté de la base carrée mesure $7$cm. Calcule le coefficient de réduction.
\item Déduis-en le volume $\mathcal{V'}$ de la pyramide sur la maquette. donne la valeur exacte en cm$^3$ puis la valeur arrondie à l'unité.
\end{enumerate}
\end{exer}
\end{document}